% Vancouver-style reference list for CDFJ (Journal of Cybersecurity, Digital
% Forensics and Jurisprudence). Per the Author Guidelines: numbered in order of
% first citation, Arabic numerals in square brackets in the text, complete
% bibliographic information with DOIs or URLs where available.
%
% Author names are given as surname followed by initials where the name order is
% unambiguous. Four authors of reference 7 are printed in full because their
% given-name and family-name order cannot be determined with confidence from the
% source, and initialising them would risk misattribution.

\begin{thebibliography}{17}

\bibitem{md_cyber_council_2025}
Maryland Cybersecurity Council. Maryland Cybersecurity Council biennial
activities report. Adelphi (MD): University of Maryland Global Campus; 2025.
Available from: \url{https://www.umgc.edu/mdcybersecuritycouncil}

\bibitem{md_local_gov_2021}
Ad Hoc Committee on State and Local Cybersecurity, Maryland Cybersecurity
Council. Maryland state and local government cybersecurity: analysis and
recommendations. Adelphi (MD): University of Maryland Global Campus; 2021.
Available from:
\url{https://www.umgc.edu/content/dam/umgc/documents/upload/maryland-state-and-local-government-cybersecurity-analysis-and-recommendations.pdf}

\bibitem{nist800_53_2020}
Joint Task Force. Security and privacy controls for information systems and
organizations. Gaithersburg (MD): National Institute of Standards and
Technology; 2020. Report No.: NIST SP 800-53 Rev. 5. Includes updates as of 10
December 2020. doi:10.6028/NIST.SP.800-53r5

\bibitem{fedramp_2022}
FedRAMP Program Management Office. FedRAMP authorization boundary guidance.
Washington (DC): U.S. General Services Administration; 2022. Rescinded and
replaced by the FedRAMP Minimum Assessment Scope standard, 2025. Available from:
\url{https://www.fedramp.gov/}

\bibitem{cmmc_2022}
Office of the Under Secretary of Defense for Acquisition and Sustainment.
Cybersecurity Maturity Model Certification (CMMC) model overview, version 2.0.
Washington (DC): U.S. Department of Defense; 2021. Superseded in part by the
CMMC Program final rule, 32 CFR Part 170, effective 16 December 2024. Available
from: \url{https://dodcio.defense.gov/CMMC/}

\bibitem{hassani_2024}
Hassani S, Sabetzadeh M, Amyot D, Liao J. Rethinking legal compliance
automation: opportunities with large language models. In: Proceedings of the
32nd IEEE International Requirements Engineering Conference (RE). Piscataway
(NJ): IEEE; 2024. p. 432--40. doi:10.1109/RE59067.2024.00051

\bibitem{policy_gap_llm_2026}
Oh Huan Lin, Jay Yong Jun Jie, Mandy Lee Ling Siu, Jonathan Pan. Automated
post-incident policy gap analysis via threat-informed evidence mapping using
large language models. arXiv:2601.03287 [preprint]. 2026. Available from:
\url{https://arxiv.org/abs/2601.03287}

\bibitem{compliancenlp_2026}
Guo D, Wu J, Yiu SM. ComplianceNLP: knowledge-graph-augmented RAG for
multi-framework regulatory gap detection. arXiv:2604.23585 [preprint]. 2026.
Available from: \url{https://arxiv.org/abs/2604.23585}

\bibitem{rethink_compliance_nlp_2025}
Kumar B, Roussinov D. NLP-based regulatory compliance: using GPT-4.0 to decode
regulatory documents. In: Proceedings of the Georg Nemetschek Institute
Symposium on Artificial Intelligence for the Built World; 2024. Available from:
\url{https://arxiv.org/abs/2412.20602}

\bibitem{md_doit_web}
Maryland Department of Information Technology. State minimum cybersecurity
standards best practices guidebook [Internet]. Annapolis (MD): Maryland
Department of Information Technology; 2023 [cited 2026 Jun]. Available from:
\url{https://doit.maryland.gov/policies/ci/Pages/state-minimum-cybersecurity-standards-best-practices-guidebook.aspx}

\bibitem{motlagh_2024}
Motlagh FN, Hajizadeh M, Majd M, Najafi P, Cheng F, Meinel C. Large language
models in cybersecurity: state-of-the-art. arXiv:2402.00891 [preprint]. 2024.
Available from: \url{https://arxiv.org/abs/2402.00891}

\bibitem{salman_2024}
Salman A, Creese S, Goldsmith M. Position paper: leveraging large language
models for cybersecurity compliance. In: Proceedings of the IEEE European
Symposium on Security and Privacy Workshops. Piscataway (NJ): IEEE; 2024.
p. 496--503.

\bibitem{xrouter_2025}
Qian C, et al. xRouter: training cost-aware LLMs orchestration system via
reinforcement learning. arXiv:2510.08439 [preprint]. 2025. Available from:
\url{https://arxiv.org/abs/2510.08439}

\bibitem{multillm_bench_2026}
Kulkarni S, Kulkarni Y. Benchmarking multi-agent LLM architectures for financial
document processing: a comparative study of orchestration patterns,
cost-accuracy tradeoffs and production scaling strategies. arXiv:2603.22651
[preprint]. 2026. Available from: \url{https://arxiv.org/abs/2603.22651}

\bibitem{pdfplumber}
Singer-Vine J. pdfplumber: plumb a PDF for detailed information about each char,
rectangle, line, and so on, and easily extract text and tables [Internet]. 2024
[cited 2026 Jun]. Available from: \url{https://github.com/jsvine/pdfplumber}

\bibitem{reportlab}
ReportLab Inc. ReportLab PDF library [Internet]. 2024 [cited 2026 Jun].
Available from: \url{https://www.reportlab.com/}

\bibitem{trivedi_ccsce_2026}
Trivedi D. Benchmarking municipal cybersecurity readiness through automated
policy analytics: evidence from Maryland local governments. J Comput Sci Coll.
2026;42(3). In press. Presented at the CCSC Eastern Conference 2026.

\end{thebibliography}
